% ============================================================
%  Paper 1 — CODA Project
%  Krylov Complexity, de Sitter Holography, and the MOND
%  Interpolating Function
%  Shane Hartley
%  JCAP submission draft — April 2026
% ============================================================
\documentclass[a4paper,11pt]{article}

% ── Packages ────────────────────────────────────────────────
\usepackage[utf8]{inputenc}
\usepackage[T1]{fontenc}
\usepackage{lmodern}
\usepackage{amsmath,amssymb,amsthm}
\usepackage{mathtools}
\usepackage{physics}          % \grad, \div, \abs, \norm
\usepackage{graphicx}
\usepackage{booktabs}
\usepackage{array}
\usepackage{xcolor}
\usepackage{hyperref}
\usepackage{cleveref}
\usepackage[numbers,sort&compress]{natbib}
\usepackage{geometry}
\usepackage{microtype}
\usepackage{setspace}
\usepackage{caption}
\usepackage[separate-uncertainty=true]{siunitx}

\geometry{
  a4paper,
  top=25mm, bottom=25mm,
  left=25mm, right=25mm
}

\hypersetup{
  colorlinks = true,
  linkcolor  = blue!70!black,
  citecolor  = green!50!black,
  urlcolor   = blue!70!black,
  pdftitle   = {Krylov Complexity, de Sitter Holography, and the MOND
                Interpolating Function},
  pdfauthor  = {Shane Hartley}
}

% ── Custom commands ──────────────────────────────────────────
\newcommand{\CK}{\mathcal{C}_K}                     % Krylov density field
\newcommand{\CKz}{\mathcal{C}_K^{(0)}}              % pure dS value
\newcommand{\CdS}{\mathcal{C}_{dS}}                  % dS complexity
\newcommand{\tausk}{\tau_*}                          % limiting surface tau
\newcommand{\Estar}{E_*}                             % limiting surface E
\newcommand{\gstar}{g_*}                             % growth rate at LES
\newcommand{\kap}{\kappa}                            % gradient constant
\newcommand{\az}{a_0}                                % Milgrom scale
\newcommand{\aN}{a_N}                                % Newtonian acceleration
\newcommand{\ab}{a_b}                                % baryonic acceleration
\newcommand{\elldS}{\ell_{\rm dS}}                   % de Sitter radius
\newcommand{\TdS}{T_{\rm dS}}                        % de Sitter temperature
\newcommand{\SdS}{S_{\rm dS}}                        % de Sitter entropy
\newcommand{\Mpl}{M_P}                               % Planck mass
\newcommand{\GN}{G_N}                                % Newton's constant
\newcommand{\dS}[1]{\mathrm{dS}_{#1}}               % de Sitter notation
\newcommand{\LES}{\mathrm{LES}}                      % limiting extremal surface
\newcommand{\AQUAL}{\mathrm{AQUAL}}                  % AQUAL action
\newcommand{\rhobar}{\rho_b}                         % baryonic density
\newcommand{\aobs}{a_{\rm obs}}                      % observed acceleration
\newcommand{\boxedeq}[1]{%
  \begin{empheq}[box=\fbox]{equation}#1\end{empheq}}

% For boxed equations without empheq
\usepackage{empheq}

% Theorem environments
\newtheorem{theorem}{Theorem}
\newtheorem{corollary}{Corollary}[theorem]

% ── Title block ──────────────────────────────────────────────
\title{\large\bfseries Krylov Complexity, de~Sitter Holography, \\[4pt]
       and the MOND Interpolating Function}

\author{Shane Hartley \\[4pt]
  \small\itshape Independent researcher \\
  \small \texttt{shane.hartley06@gmail.com}}

\date{April 2026 \quad --- \quad Preprint, submitted to JCAP}

% ============================================================
\begin{document}
\maketitle

% ── Abstract ────────────────────────────────────────────────
\begin{abstract}
We show that the MOND interpolating function $\mu(x)=\tanh(x)$ and the
Milgrom acceleration scale $\az = cH_0/2\pi$ can be derived from
holographic Krylov complexity in de~Sitter spacetime, using the DSSYK
correspondence established by Heller et~al.\ \cite{Heller:2025} and
Aguilar-Gutierrez \cite{Aguilar:2024}.  The derivation has four steps.
First, we compute the Krylov complexity density of the DSSYK limiting
extremal surface in $\dS{4}$, obtaining
$\CKz = 2/(3\sqrt{3}\,\GN\ell^2)$; this satisfies all covariance
requirements and is state-independent.  Second, we show that the double
degeneracy of the limiting surface forces a non-analytic response to
Newtonian perturbations --- $(\delta\tausk)^2=-3\Phi/4$ --- yielding the
gradient relation $|\nabla\CK|=\kap\,\aN$.  Third, differentiating the
DSSYK classical complexity function $C_K(x)=\log\cosh(x)$ yields
$\mu(x)=\tanh(x)$ with the correct deep-MOND and Newtonian asymptotes.
Fourth, the coupling $\kap$ cancels identically when the kinetic action
is written in AQUAL form, producing
\begin{equation*}
  \nabla\cdot\!\left[\tanh\!\left(\frac{\aN}{\az}\right)\nabla\Phi\right]
  = 4\pi G\rhobar
\end{equation*}
with $\az = cH_0/2\pi$ from the de~Sitter temperature --- within 10\%
of the observed Milgrom scale with no free parameters.  We test the
predicted interpolating function against all 175 SPARC galaxies and find
fits competitive with the standard simple function, with $\mu=\tanh$
performing best in the deep-MOND dwarf-galaxy regime.  A precisely
identified open problem --- deriving the AQUAL form of the kinetic action
from the DSSYK second variation --- is the remaining step to complete the
first-principles derivation.
\end{abstract}

\noindent\textbf{Keywords:} MOND $\cdot$ Krylov complexity $\cdot$
holography $\cdot$ de~Sitter $\cdot$ DSSYK $\cdot$ galactic dynamics

% ============================================================
\section{Introduction}
\label{sec:intro}

\subsection{The MOND problem}

Modified Newtonian Dynamics (MOND) \cite{Milgrom:1983a,Milgrom:1983b} is
the most successful empirical description of galactic dynamics without
particle dark matter.  The single relation
\begin{equation}
  \aN\,\mu\!\left(\frac{\aN}{\az}\right) = \ab,
  \qquad \az \simeq 1.2\times10^{-10}\,\si{m.s^{-2}},
  \label{eq:MOND}
\end{equation}
where $\ab$ is the Newtonian acceleration from baryons alone, accounts for
the flat rotation curves of all 175 galaxies in the SPARC catalogue
\cite{Lelli:2016}, predicts the baryonic Tully-Fisher relation
$v_c^4=\az G M_b$ \cite{Verheijen:2001,McGaugh:2005}, and encodes the
Radial Acceleration Relation (RAR) with scatter of only 0.13\,dex over
five decades in acceleration \cite{McGaugh:2016}.  Two quantities drive
this success: the acceleration scale $\az$ and the interpolating function
$\mu(x)$.

Both remain without a derivation from fundamental physics.  The scale
$\az\approx cH_0/6$ is numerologically close to $cH_0$, suggesting a
cosmological origin.  Several frameworks have connected $\az$ to the
Hubble constant through thermodynamic arguments
\cite{Milgrom:2023,Zhao:2010,Verlinde:2017}, but these do not determine
$\mu$.  The function $\mu$ is postulated in every existing relativistic
completion of MOND \cite{Skordis:2021,Bekenstein:2004,Deffayet:2011} ---
it is fit to data or adopted on phenomenological grounds.  The question of
\emph{why} $\mu$ takes a specific form, and \emph{why} the transition from
Newtonian to MOND dynamics occurs precisely at $\aN\sim cH_0$, has no
answer in the current literature.

\subsection{Holographic Krylov complexity as the new ingredient}

Krylov complexity \cite{Parker:2019,Avdoshkin:2022} is an operator-theoretic
measure of quantum information spreading.  For an operator $O(t)$ evolving
under Hamiltonian $H$, the Krylov complexity $K(t)$ quantifies how far
$O(t)$ has spread in the Krylov basis generated by repeated commutators
with $H$.  For maximally chaotic systems it grows as $K(t)\propto
e^{\lambda_K t}$, saturating the MSS chaos bound with Lyapunov exponent
$\lambda_K = 2\pi T$.

Recent work by Heller et~al.\ \cite{Heller:2025} and Aguilar-Gutierrez
\cite{Aguilar:2024} establishes the first microscopic computation of
Krylov complexity in a cosmological spacetime.  In the
DSSYK/LdS$_2$/SdS$_3$ holographic triality, the Krylov complexity of the
Hartle-Hawking state equals the volume of a timelike extremal slice in de
Sitter, and the classical complexity function is
\begin{equation}
  C_K(x) = \log\cosh(x).
  \label{eq:CK}
\end{equation}
This function has a natural derivative: $dC_K/dx = \tanh(x)$.  The
central observation of this paper is that this derivative \emph{is} the
MOND interpolating function.  The identification follows from a bulk
calculation showing that the gradient of the holographic complexity
density equals $\kap\aN$ for a constant $\kap$ that cancels in the MOND
action.

\subsection{Scope of this work}
\label{sec:scope}

We derive $\mu(x)=\tanh(x)$ and $\az=cH_0/2\pi$ from the DSSYK
holographic construction, given that the kinetic action coupling $\CK$ to
gravity is of AQUAL form.  The AQUAL form is motivated by a partial
computation of the DSSYK second variation (\cref{sec:second_var}), which
shows the correct non-analytic power structure, but is not yet fully
derived from the holographic bulk.  This open step --- the holographic
dictionary problem --- is precisely identified in \cref{sec:open}.

This is the first derivation of a specific MOND interpolating function
from an underlying quantum-gravitational framework.  The function
$\mu=\tanh$ was previously known to work phenomenologically
\cite{Zhao:2006}, but had no theoretical motivation.  The approach is
complementary to entropy-based derivations of $\az$
\cite{Milgrom:2023,Zhao:2010,Verlinde:2017} in that it produces both $\az$
and $\mu$ simultaneously from the same DSSYK complexity function.

\subsection{Organisation}

\Cref{sec:DSSYK} reviews the DSSYK holographic complexity programme.
\Cref{sec:density} computes the explicit Krylov density in $\dS{4}$.
\Cref{sec:perturbed} derives the gradient relation.
\Cref{sec:MOND} presents the MOND equation.
\Cref{sec:SPARC} compares predictions against SPARC.
\Cref{sec:discussion} discusses connections to existing frameworks.
\Cref{sec:conclusions} concludes.

\medskip
\noindent\textbf{Conventions.}  Signature $(-,+,+,+)$; $c$ restored where
relevant; $\GN$ Newton's constant; $\ell=c/H_0$ de~Sitter radius;
$\Mpl^2=(8\pi\GN)^{-1}$.

% ============================================================
\section{DSSYK Holographic Krylov Complexity in de~Sitter}
\label{sec:DSSYK}

\subsection{The DSSYK model and its de~Sitter limit}

The Double-Scaled Sachdev-Ye-Kitaev (DSSYK) model is a disordered
quantum-mechanical system --- a large-$N$ limit of the SYK model with
$p\sim\sqrt{N}$ fermion interactions --- which flows in a particular limit
to two-dimensional de~Sitter gravity (sine-dilaton gravity)
\cite{Aguilar:2024,Blommaert:2024,Susskind:2022}.  The relevant parameter
is the chord length $\theta\in(0,\pi)$, with the de~Sitter limit
$\theta\to\pi$.

In this limit, the DSSYK model is dual to quantum gravity in a static patch
of $\mathrm{dS}_2$ (LdS$_2$), which is related to a static patch of
$\mathrm{dS}_3$ (SdS$_3$) by dimensional oxidation.  The Hartle-Hawking
no-boundary state corresponds to a specific state in the DSSYK chord
Hilbert space.  Heller et~al.\ \cite{Heller:2025} prove that the Krylov
spread complexity of this state equals the volume of a timelike extremal
codimension-1 slice in the bulk $\mathrm{dS}_2$:
\begin{equation}
  2|\log q|\,C_K(\chi)_{\theta\approx\pi} = -i\,L_{\rm dS}(\chi),
  \label{eq:C=V}
\end{equation}
where $q$ is the DSSYK deformation parameter and $L_{\rm dS}(\chi)$ is
the (imaginary) volume of the extremal slice anchored at boundary time
$\chi/2$.  This is the de~Sitter analogue of the AdS Complexity=Volume
conjecture \cite{Brown:2016}.

\subsection{The classical complexity function}

The classical saddle-point result \cite[Eq.~20]{Heller:2025} is
\begin{equation}
  C_K(\chi) \propto \log\!\left[\cosh\!\left(\frac{\chi\theta}{2}\right)\right],
  \label{eq:CK_classical}
\end{equation}
exhibiting a characteristic quadratic-to-linear transition:
\begin{itemize}
  \item \textbf{Precursor regime} ($\chi\ll\elldS$):
    $C_K\approx(\chi\theta/2)^2/2$ --- quadratic growth.
  \item \textbf{Scrambling regime} ($\chi\gg\elldS$):
    $C_K\approx\chi\theta/2 - \log 2$ --- linear growth.
\end{itemize}
The transition occurs at $\chi_*\sim\elldS=c/H_0$.
The late-time growth rate \cite[Eq.~70]{Heller:2025} is
\begin{equation}
  \frac{dC_K}{d\chi}\bigg|_{\rm late}
  \propto \SdS\cdot\TdS
  = \frac{\pi\ell^2}{\GN}\cdot\frac{H_0}{2\pi},
  \label{eq:growth_rate}
\end{equation}
where $\SdS=\pi\ell^2/\GN$ is the Gibbons-Hawking entropy and
$\TdS=H_0/2\pi$ is the de~Sitter temperature.

The key property of $C_K(x)=\log\cosh(x)$ relevant to MOND is
\begin{equation}
  \frac{dC_K}{dx} = \tanh(x).
  \label{eq:mu_from_CK}
\end{equation}
In \cref{sec:MOND} we show this becomes the MOND interpolating function
at $x=\aN/\az$.

\subsection{The bulk extremal surface}

Heller et~al.\ propose the holographic complexity in $\mathrm{dS}_{d+1}$
\cite[Eq.~23]{Heller:2025}:
\begin{equation}
  \CdS \equiv \frac{-i\,V}{\GN\,\ell},
  \label{eq:CdS_def}
\end{equation}
where $V$ is the volume of the timelike extremal codimension-1 slice.
In coordinates $(\tau,w,\Omega_{d-1})$ with metric
\begin{equation}
  ds^2 = \frac{\ell^2}{\tau^2}\!\left[
    -\frac{d\tau^2}{1-\tau^2} + (1-\tau^2)\,dw^2 + d\Omega_{d-1}^2
  \right],
  \label{eq:dS_metric}
\end{equation}
the extremal slice volume is
\begin{equation}
  V(w_0) = \ell^d\,\Omega_{d-1}\int d\tau\;
  \frac{i}{\tau^d}\sqrt{\frac{1}{1-\tau^2}-(1-\tau^2)\dot{w}^2(\tau)},
  \label{eq:V_integral}
\end{equation}
with $\dot{w}=dw/d\tau$.  As $w_0\to\infty$, embeddings accumulate at the
\textbf{limiting extremal surface} (LES), whose turning point $\tausk$
depends only on the de~Sitter geometry --- not on the quantum state.  The
LES is the late-time attractor of the complexity slices; its position is
determined by stationarity of the volume growth rate $dV/dw_0$
(\cref{sec:density}).

\subsection{Three routes to the MOND scale}
\label{sec:three_routes}

The DSSYK programme identifies $\az\sim cH_0$ through three independent
routes:

\begin{enumerate}
  \item \textbf{Chaos bound \cite{Aguilar:2024}:}
    The Lyapunov exponent saturates the MSS bound, $\lambda_K=H_0$, giving
    $\tau_*=H_0^{-1}$ and $\az=c/\tau_*=cH_0$.

  \item \textbf{Complexity transition \cite{Heller:2025}:}
    The quadratic-to-linear transition in $C_K$ at
    $\chi_*\sim\elldS$ gives $\az=c^2/\elldS=cH_0$.

  \item \textbf{De~Sitter temperature \cite{Heller:2025}:}
    The growth rate $\propto\SdS\cdot\TdS$ identifies $\az = c\TdS =
    cH_0/2\pi$.
\end{enumerate}

Routes~1--2 give $cH_0$; Route~3 gives $cH_0/2\pi$ from the thermal
normalisation.  Route~3 is observationally preferred:
$cH_0/2\pi\approx1.082\times10^{-10}\,\si{m.s^{-2}}$ versus observed
$\az^{\rm obs}\approx1.2\times10^{-10}\,\si{m.s^{-2}}$ (10\% agreement,
no free parameters).  For $H_0=75\,\si{km.s^{-1}.Mpc^{-1}}$ (SH0ES
\cite{Riess:2022}) this tightens to 3\%.

% ============================================================
\section{The Krylov Complexity Density in $\dS{4}$}
\label{sec:density}

\subsection{Setup}

We work in $d=3$ ($\dS{4}$).  The extremal slice energy $E$ and turning
point $\tausk$ satisfy
\begin{equation}
  E^2 = \frac{\tausk^2-1}{\tausk^{2d}}, \qquad
  f(\tausk,E) \equiv 1-\tausk^2+E^2\tausk^{2d} = 0,
  \label{eq:constraint}
\end{equation}
and the volume growth rate is
\begin{equation}
  g(\tausk) \equiv \tausk^{-d}\sqrt{\tausk^2-1}.
  \label{eq:growth}
\end{equation}
The late-time volume satisfies $dV/dw_0\propto g(\tausk)$.

\subsection{The limiting surface}

As $w_0\to\infty$, $\tausk$ accumulates to the stationary point
$g'(\tausk)=0$.  For $d=3$:
\begin{equation}
  g'(\tausk)
  = \frac{\tausk^{-4}(3-2\tausk^2)}{\sqrt{\tausk^2-1}}\bigg|_*
  = 0 \implies \tausk^2 = \frac{3}{2},
\end{equation}
\begin{empheq}[box=\fbox]{equation}
  \tausk^{(d=3)} = \sqrt{\tfrac{3}{2}} \approx 1.2247.
  \label{eq:tau_star}
\end{empheq}
The corresponding energy and growth-rate coefficient are
\begin{equation}
  \Estar = g(\tausk) = \frac{2}{3\sqrt{3}},
  \qquad
  \frac{dV}{dw_0}\bigg|_* = \frac{8\pi i\ell^3}{3\sqrt{3}},
  \qquad
  c_{d=3} = \frac{16\pi}{3\sqrt{3}} \approx 9.68,
  \label{eq:E_star}
\end{equation}
in agreement with \cite[Fig.~3]{Heller:2025} for $d=3$.

\subsection{Double degeneracy of the limiting surface}
\label{sec:double_deg}

A structural feature critical for \cref{sec:perturbed}: both partial
derivatives of the constraint function vanish simultaneously at
$(\tausk,\Estar)$:
\begin{equation}
  \frac{\partial f}{\partial\tau}\bigg|_* = 0
  \quad\text{(turning-point definition)}, \qquad
  g'(\tausk) = 0
  \quad\text{(LES definition)}.
  \label{eq:double_deg}
\end{equation}
These are generically independent conditions; their simultaneous vanishing
makes $(\tausk,\Estar)$ a \emph{degenerate critical point} of the extremal
surface family.  The consequence is that metric perturbations shift
$\tausk$ by $|\Phi|^{1/2}$, not $|\Phi|$ --- first-order perturbation
theory breaks down, forcing a non-analytic response.  This geometric
non-analyticity is the origin of MOND non-linearity
(\cref{sec:perturbed}).

\subsection{The Krylov complexity density}
\label{sec:CK_density}

Normalising the late-time growth rate to a 4D proper volume from the Heller
et~al.\ construction in $d=3$:
\begin{empheq}[box=\fbox]{equation}
  \CKz = \frac{2}{3\sqrt{3}\,\GN\ell^2}
  \approx 0.385\,\Mpl^2 H_0^2.
  \label{eq:CK0}
\end{empheq}
This satisfies the four requirements for inclusion in a covariant action:

\begin{table}[h]
\centering
\renewcommand{\arraystretch}{1.3}
\begin{tabular}{lll}
  \toprule
  Requirement & Status & Comment \\
  \midrule
  Scalar under diffeomorphisms & $\checkmark$ &
    Metric-determined geometric volume \\
  Dimensions $[L^{-4}]$ & $\checkmark$ &
    $[\GN\ell^2]^{-1}\sim\Mpl^2 H_0^2$ \\
  Finite & $\checkmark$ &
    Growth rate $\Estar$ finite at LES \\
  State-independent & $\checkmark$ &
    $\tausk=\sqrt{3/2}$ determined by metric alone \\
  \bottomrule
\end{tabular}
\caption{Properties of the dS$_4$ Krylov complexity density $\CKz$.}
\label{tab:checks}
\end{table}

In pure de~Sitter, $\CK$ is spatially constant (correct by maximal
symmetry).  Spatial variation arises only from perturbations of the dS
background, computed next.

% ============================================================
\section{Perturbed de~Sitter --- the Gradient Relation}
\label{sec:perturbed}

\subsection{Setup}

Consider a conformal perturbation in Newtonian gauge:
\begin{equation}
  g_{\mu\nu} \to (1+2\Phi)\,g_{\mu\nu}^{(0)}, \qquad |\Phi|\ll 1,
  \label{eq:Newtonian_gauge}
\end{equation}
where $\Phi(x)$ is the Newtonian gravitational potential.  This shifts the
LES turning point $\tausk\to\tausk+\delta\tau(x)$.

\subsection{Non-analytic response from the double degeneracy}

The perturbed constraint $f_{\rm pert}(\tausk+\delta\tau, \Estar+\delta
E, \Phi)=0$, expanded around pure dS.  Using $\partial
f/\partial\tau|_*=0$ (double degeneracy), the leading term is quadratic
in $\delta\tau$:
\begin{equation}
  \frac{1}{2}f_{\tau\tau}\big|_*\,(\delta\tau)^2 + K\,\Phi
  = 0 + O(\delta\tau^3,\Phi^2),
\end{equation}
where $f_{\tau\tau}|_*=8$ and $K=3$ (computed at $\tausk=\sqrt{3/2}$,
$\Estar=2/(3\sqrt{3})$, using $\Estar^2\tausk^6=1/2$).  Solving:
\begin{empheq}[box=\fbox]{equation}
  (\delta\tau)^2 = -\frac{3\Phi}{4}.
  \label{eq:delta_tau}
\end{empheq}
For attractive gravity ($\Phi<0$), $\delta\tau\in\mathbb{R}$.  The
turning point shifts as $\delta\tau\propto|\Phi|^{1/2}$ --- the
double-degeneracy non-analyticity.

\subsection{Complexity density response}

Since $g'(\tausk)=0$, the leading change in growth rate is
\begin{equation}
  \delta g = \frac{1}{2}g''(\tausk)(\delta\tau)^2 + O(|\Phi|^{3/2}),
\end{equation}
with $g''(\tausk)=-16\sqrt{3}/9$ and $\gstar=\Estar=2/(3\sqrt{3})$:
\begin{equation}
  \frac{\delta g}{\gstar}
  = \frac{\tfrac{1}{2}\cdot(-\tfrac{16\sqrt{3}}{9})\cdot(-\tfrac{3\Phi}{4})}
         {\tfrac{2\sqrt{3}}{9}}
  = 3\Phi.
\end{equation}
\begin{empheq}[box=\fbox]{equation}
  \frac{\delta\CK(x)}{\CKz} = 3\,\Phi(x).
  \label{eq:dCK}
\end{empheq}
The complexity density tracks the Newtonian potential with coefficient 3
--- a pure number from the $\dS{4}$ LES geometry.

\subsection{The gradient relation}

Taking the spatial gradient of \cref{eq:dCK}:
\begin{equation}
  \nabla\CK(x) = 3\CKz\,\nabla\Phi(x) \equiv \kap\,\nabla\Phi(x),
  \label{eq:grad_CK}
\end{equation}
where
\begin{equation}
  \kap \equiv 3\CKz = \frac{2}{\sqrt{3}\,\GN\ell^2}.
\end{equation}
Since $\aN\equiv|\nabla\Phi|$:
\begin{empheq}[box=\fbox]{equation}
  |\nabla\CK(x)| = \kap\,\aN(x).
  \label{eq:grad_relation}
\end{empheq}
\textbf{The gradient of the Krylov complexity density equals $\kap$ times
the local Newtonian acceleration}, with $\kap$ fixed by the $\dS{4}$
geometry.  As shown in \cref{sec:kappa_cancel}, $\kap$ cancels identically
in the MOND action.

% ============================================================
\section{The CODA MOND Equation}
\label{sec:MOND}

\subsection{Motivation for the AQUAL kinetic action}
\label{sec:AQUAL_motivation}

We couple $\CK$ to gravity through a kinetic action.  Two requirements
fix the form: (i) the MOND equation must be second order in $\Phi$,
ruling out linear couplings $\propto\CK$; (ii) the action must reduce
to the Poisson equation for $\aN\gg\az$.  These are satisfied by an
action of AQUAL form \cite{Bekenstein:1984}:
\begin{equation}
  S_{\rm kin} = -\frac{\az^2}{8\pi G}\int d^4x\sqrt{-g}\;
  F\!\left(\frac{|\nabla\CK|^2}{\kap^2\,\az^2}\right),
  \label{eq:Skin}
\end{equation}
with $F(s)\to s$ for $s\gg1$.  Additional motivation comes from the
second variation $\delta^2\CdS/\delta g^2$, which produces a non-analytic
$|\Phi|^{3/2}$ term with the correct power for deep-MOND AQUAL
$F(s)\sim s^{3/2}$ (\cref{sec:second_var}).  We adopt the AQUAL form as
a working hypothesis; deriving it from the DSSYK bulk is the primary open
problem (\cref{sec:open}).

The full CODA action is
\begin{equation}
  S_{\rm CODA} = \int d^4x\sqrt{-g}\!\left[
    \frac{R}{16\pi G}
    + \xi\,C_{\alpha\beta\gamma\delta}C^{\alpha\beta\gamma\delta}
    - \frac{\az^2}{8\pi G}\,F\!\left(\frac{|\nabla\CK|^2}{\kap^2\az^2}\right)
  \right] + S_m,
  \label{eq:Stotal}
\end{equation}
where the Weyl-squared term (with ghost mass $m_2^2=1/(32\pi G\xi)$
\cite{Stelle:1977}) is Bach-flat in all Einstein geometries and does not
affect galactic dynamics.

\subsection{The $\kappa$ cancellation}
\label{sec:kappa_cancel}

Substituting $|\nabla\CK|^2=\kap^2\aN^2$ from \cref{eq:grad_relation}:
\begin{equation}
  \frac{|\nabla\CK|^2}{\kap^2\,\az^2}
  = \frac{\kap^2\aN^2}{\kap^2\az^2}
  = \frac{\aN^2}{\az^2}.
  \label{eq:kappa_cancel}
\end{equation}
The coupling $\kap$ cancels identically, and $S_{\rm kin}$ reduces to
\begin{equation}
  S_{\rm kin} = -\frac{\az^2}{8\pi G}\int d^4x\sqrt{-g}\;
  F\!\left(\frac{\aN^2}{\az^2}\right),
  \label{eq:Skin_reduced}
\end{equation}
exactly the AQUAL action of \citeauthor{Bekenstein:1984}
\cite{Bekenstein:1984}.  The cancellation is robust: any
$|\nabla\CK|=\kap'\aN$ produces the same result.  The MOND equation is
therefore independent of the holographic normalisation.

\subsection{The MOND interpolating function from DSSYK}
\label{sec:mu_from_DSSYK}

The AQUAL function $F$ determines $\mu$ via $F'(s)=\mu(\sqrt{s})$.
Identifying $F$ with the DSSYK complexity:
\begin{equation}
  \mu(x) = \frac{dC_K}{dx}\bigg|_{C_K=\log\cosh} = \tanh(x),
  \label{eq:mu_tanh}
\end{equation}
evaluated at $x=\aN/\az$.  The corresponding AQUAL function is
\begin{equation}
  F(s) = 2\sqrt{s}\,\log\cosh(\sqrt{s})
       - 2\int_0^{\sqrt{s}}\log\cosh(u)\,du,
  \label{eq:F}
\end{equation}
which is an integral transform of $C_K$: $F(s)=2\sqrt{s}\,C_K(\sqrt{s})
-2\int_0^{\sqrt{s}}C_K(u)\,du$.  The limits are:
\begin{align}
  s\ll 1: &\quad F(s)\to\tfrac{2}{3}s^{3/2} \implies \mu(x)\to x
    \quad\text{(deep MOND)}\ \checkmark \label{eq:deep_MOND}\\
  s\gg 1: &\quad F(s)\to s \implies \mu(x)\to 1
    \quad\text{(Newtonian)}\ \checkmark \label{eq:Newton}
\end{align}
Both limits are verified analytically and numerically.
See \cref{fig:mu} for a comparison with existing interpolating functions.

\subsection{The MOND theorem and corollaries}
\label{sec:theorem}

\begin{theorem}[CODA MOND Theorem]
  Given the CODA kinetic action \eqref{eq:Skin} with $F'(s)=\tanh(\sqrt{s})$
  and the gradient relation \eqref{eq:grad_relation}, the quasi-static
  weak-field field equation is
  \begin{empheq}[box=\fbox]{equation}
    \nabla\cdot\!\left[\tanh\!\left(\frac{\aN}{\az}\right)\nabla\Phi\right]
    = 4\pi G\,\rhobar,
    \label{eq:MOND_eq}
  \end{empheq}
  with $\az = c\TdS = cH_0/2\pi$.
\end{theorem}

\noindent The proof has four steps: (i) $\delta\CK=3\Phi\CKz$ from
\eqref{eq:dCK}; (ii) $|\nabla\CK|=\kap\aN$ from \eqref{eq:grad_relation};
(iii) $\kap$ cancels identically \eqref{eq:kappa_cancel}; (iv)
$F'(s)=\tanh(\sqrt{s})$ from \eqref{eq:F}.

\begin{corollary}[Baryonic Tully-Fisher Relation]
  In the deep-MOND limit $\aN\ll\az$, $\tanh(\aN/\az)\approx\aN/\az$
  and the asymptotic circular velocity satisfies
  \begin{equation}
    v_c^4 = \az\,G\,M_b.
    \label{eq:BTFR}
  \end{equation}
  Using $\az=cH_0/2\pi$: $v_c=109\,\si{km/s}$ for $M_b=10^{10}M_\odot$
  and $v_c=195\,\si{km/s}$ for $M_b=10^{11}M_\odot$, consistent with the
  observed BTFR zero-point \cite{Verheijen:2001,McGaugh:2005}.
\end{corollary}

\begin{corollary}[Radial Acceleration Relation]
  The equation $\aN\tanh(\aN/\az)=\ab$ gives a unique $\aN$ for each
  $\ab$, with $\aN\to\ab$ for $\ab\gg\az$ and $\aN\to\sqrt{\ab\az}$ for
  $\ab\ll\az$ --- a parameter-free prediction of the observed RAR
  \cite{McGaugh:2016}.  See \cref{fig:RAR}.
\end{corollary}

\subsection{What is proved and what remains open}
\label{sec:open}

\paragraph{What is proved.}
Given the AQUAL kinetic action (premise~iii), the MOND equation
\eqref{eq:MOND_eq} follows exactly with $\az=cH_0/2\pi$.

\paragraph{What remains open.}
Premise~(iii) --- the AQUAL form of the kinetic action --- is supported by
the second variation (\cref{sec:second_var}) but not yet derived from the
DSSYK bulk.  The required calculation is to show:
\begin{equation}
  \frac{\delta^2\CdS}{\delta\Phi(x)\,\delta\Phi(y)}
  \xrightarrow[\text{dict.}]{\text{holo.}}
  -\frac{\az^2}{8\pi G}\,F''\!\left(\frac{\aN^2}{\az^2}\right)
  \cdot\frac{\partial\aN^2}{\partial\Phi(x)}\,\frac{\partial\aN^2}{\partial\Phi(y)},
  \label{eq:open}
\end{equation}
i.e.\ that the DSSYK boundary time responds to $|\nabla\Phi|$ (gradient)
rather than $|\Phi|$ (value).  The holographic dictionary problem is
analysed in detail in companion work \cite{CODA:dict}.

\subsection{Evidence from the second variation}
\label{sec:second_var}

To support premise~(iii), we partially compute
$\delta^2\CdS/\delta g^2$.

\paragraph{Jacobi spectrum.}
The Jacobi operator governing shape deformations of the LES has
eigenvalues
\begin{equation}
  \lambda(k,\ell) = 3k^2 - \tfrac{3}{2}\ell(\ell+1) + \tfrac{7}{3}
  \label{eq:Jacobi}
\end{equation}
for mode $(k,\ell)$.  The ground state $\lambda_0=7/3\neq0$: there is no
Jacobi near-zero mode, and MOND does not arise from a shape instability
of the LES.

\paragraph{Non-analytic MOND term.}
The MOND signal arises from the triple derivative
$g'''(\tausk)=208\sqrt{2}/9$ acting through the double degeneracy.
The third-order expansion of the growth rate gives
\begin{equation}
  g = \Estar\!\left[1 + 3\Phi
    + \frac{13\sqrt{2}}{2}\,|\Phi|^{3/2} + O(\Phi^2)\right].
  \label{eq:g_expand}
\end{equation}
The non-analytic $|\Phi|^{3/2}$ term matches the deep-MOND AQUAL limit
$F(s)\sim s^{3/2}$ when $|\Phi|\sim(\aN/\az)^2$, which holds at the
MOND radius $r_M=\sqrt{GM_b/\az}$.

% ============================================================
\section{Comparison with SPARC Rotation Curves}
\label{sec:SPARC}

\subsection{The prediction}

With $\mu=\tanh$ and $\az=cH_0/2\pi$, we solve $\aN\tanh(\aN/\az)=\ab(r)$
at each radius, with
\begin{equation}
  \ab(r) = \frac{V_{\rm gas}|V_{\rm gas}|+\Upsilon_d V_{\rm disk}^2
              +\Upsilon_b V_{\rm bul}^2}{r},
  \label{eq:ab}
\end{equation}
from the SPARC mass models \cite{Lelli:2016}, with canonical stellar
mass-to-light ratios $\Upsilon_d=0.5\,M_\odot/L_\odot$ and
$\Upsilon_b=0.7\,M_\odot/L_\odot$ \cite{Schombert:2020}.  The
sign-preserving convention $V|V|/r$ handles the 361 negative
$V_{\rm gas}$ entries.  No parameters are fitted; the circular velocity
follows from $v_c=\sqrt{\aN r}$.

\subsection{Results across 175 galaxies}

We evaluate three models on all 175 SPARC galaxies (3391 data points):
CODA ($\mu=\tanh$, $\az=cH_0/2\pi$), tanh ($\mu=\tanh$, $\az=\az^{\rm obs}$),
and simple ($\mu=x/\sqrt{1+x^2}$ \cite{Zhao:2006}, $\az=\az^{\rm obs}$).
Results are shown in \cref{tab:sparc}.

\begin{table}[h]
\centering
\renewcommand{\arraystretch}{1.3}
\begin{tabular}{lcccc}
  \toprule
  Model & Mean $\chi^2_r$ & Median $\chi^2_r$ & RMS (km/s) & $\chi^2_r<5$ \\
  \midrule
  CODA ($\az=cH_0/2\pi$) & 40.6 & 11.2 & 18.5 & 58/175 \\
  tanh ($\az=\az^{\rm obs}$) & 36.9 & 10.6 & 18.0 & 56/175 \\
  simple ($\az=\az^{\rm obs}$) & 34.2 & 9.8 & 17.5 & 59/175 \\
  \bottomrule
\end{tabular}
\caption{Goodness-of-fit statistics for 175 SPARC galaxies with fixed
  canonical mass-to-light ratios $\Upsilon_d=0.5$, $\Upsilon_b=0.7$.
  Elevated $\chi^2_r$ values relative to fitted-M/L MOND analyses
  \cite{Famaey:2012} reflect the intentional absence of per-galaxy
  stellar mass optimisation; the relevant comparison is between columns.}
\label{tab:sparc}
\end{table}

\begin{table}[h]
\centering
\renewcommand{\arraystretch}{1.3}
\begin{tabular}{lccc}
  \toprule
  Hubble type & $n$ & CODA & tanh obs / simple obs \\
  \midrule
  Early ($T\leq3$) & 28 & 49.0 & 45.5 / 34.3 \\
  Late ($T=4$--$7$) & 66 & 11.8 & 10.1 / 9.0 \\
  Dwarf ($T\geq8$) & 81 & \textbf{7.69} & 7.69 / \textbf{7.88} \\
  \bottomrule
\end{tabular}
\caption{Median $\chi^2_r$ stratified by Hubble type.  In the deep-MOND
  dwarf regime ($T\geq8$), $\mu=\tanh$ narrowly outperforms the simple
  function (7.69 vs.\ 7.88).  The early-type deficit for the CODA model
  reflects the 10\% discrepancy in $\az$ that most strongly affects
  high-acceleration early-type galaxies near the transition.}
\label{tab:hubble}
\end{table}

\subsection{The discriminant at Euclid precision}

The two interpolating functions differ at the transition $\aN=\az$:
\begin{equation}
  \mu_{\tanh}(1) = \tanh(1) \approx 0.762, \qquad
  \mu_{\rm simple}(1) = 1/\sqrt{2} \approx 0.707 \quad (+7.7\%).
  \label{eq:mu_transition}
\end{equation}
Individual rotation curves cannot resolve this difference at typical 20\%
M/L uncertainty.  Euclid weak-lensing stacking in narrow $\ab/\az$ bins
near the transition will reach $\sim5\%$ precision \cite{Amendola:2018,
Desmond:2020}, making the discrimination accessible.

% ============================================================
\section{Discussion}
\label{sec:discussion}

\subsection{Relation to existing MOND frameworks}

\begin{table}[h]
\centering
\renewcommand{\arraystretch}{1.3}
\begin{tabular}{llcc}
  \toprule
  Function & $\mu(x)$ & $\mu(1)$ & Source \\
  \midrule
  CODA (this work) & $\tanh(x)$ & 0.762 &
    Derived from $dC_K/dx$ \\
  Simple & $x/\sqrt{1+x^2}$ & 0.707 &
    Phenomenological \cite{Zhao:2006} \\
  Standard & $x/(1+x)$ & 0.500 &
    Original Milgrom \cite{Milgrom:1983a} \\
  \bottomrule
\end{tabular}
\caption{Comparison of MOND interpolating functions.}
\label{tab:mu}
\end{table}

The tanh function was previously used on phenomenological grounds
\cite{Zhao:2006,Gentile:2011}; this paper provides the first derivation
from a theoretical framework --- specifically as the derivative of the
DSSYK Krylov complexity function.

Among relativistic completions, AeST \cite{Skordis:2021} is the closest
structural analogue.  AeST postulates a kinetic Lagrangian
$\mathcal{F}(K_{\mu\nu\alpha\beta})$ for a vector field $A^\mu$ that
reduces to AQUAL in the quasi-static limit.  CODA produces the same
reduction with $F$ derived from $C_K=\log\cosh$ rather than postulated,
and without introducing a fundamental vector field.  However, AeST's
$A^\mu$ ensures lensing-dynamics equivalence by construction, while
CODA's current framework does not guarantee this.  Constructing a CODA
vector sector is an important open problem.

Verlinde's emergent gravity \cite{Verlinde:2017} is the closest conceptual
predecessor, deriving $\az\sim cH_0$ from the entropic response of de
Sitter to mass concentrations.  CODA uses the Krylov complexity gradient
instead of Verlinde's entropic displacement, and additionally determines
$\mu$.  Whether $|\nabla\CK|$ and the entropic displacement are related
is a natural direction for future work.

\subsection{The $2\pi$ factor and the holographic dictionary}
\label{sec:2pi}

Routes~1--2 of \cref{sec:three_routes} give $\az=cH_0$; Route~3 gives
$cH_0/2\pi$.  The discrepancy is the thermal normalisation of the
de~Sitter temperature.  Resolving which route is correct requires the
holographic dictionary: if the DSSYK boundary time responds to
$|\nabla\Phi|$ (gradient), the thermal scale $\TdS=H_0/2\pi$ enters and
Route~3 is selected.  If it responds to $|\Phi|$ (value), Routes~1--2
give $cH_0$.  The non-analytic $|\Phi|^{3/2}$ term in \cref{eq:g_expand}
provides evidence for gradient response --- and for Route~3 --- but a
global proof requires the full bulk-to-boundary map.

\subsection{Galaxy clusters and gravitational lensing}

MOND without additional dark matter fails for galaxy clusters
\cite{Angus:2006}.  The present work inherits this problem from MOND
phenomenology.  A relativistic CODA framework with a vector sector may
partially address it, analogously to AeST \cite{Skordis:2021}.

\subsection{CMB and large-scale structure}

The Weyl-squared term in \eqref{eq:Stotal} is Bach-flat in all Einstein
geometries, giving corrections $(k/\Mpl)^2\sim10^{-114}$ at CMB scales.
The MOND kinetic sector is suppressed at high redshift ($\aN\gg\az$
everywhere before galaxy formation) and activates only in the
low-acceleration regime of present-day galactic disks.  A full
perturbation theory of the MOND sector is deferred to future work.

% ============================================================
\section{Conclusions}
\label{sec:conclusions}

The MOND interpolating function has been unknown theoretically for four
decades.  We have shown that it can be derived from holographic Krylov
complexity in de~Sitter spacetime.  The derivation produces
$\mu(x)=\tanh(x)$ as the derivative of the DSSYK classical complexity
function $C_K(x)=\log\cosh(x)$, and the Milgrom scale
$\az=cH_0/2\pi$ from the de~Sitter temperature --- both without free
parameters.

Three results make the derivation concrete:

\begin{enumerate}
  \item \textbf{The Krylov complexity density in $\dS{4}$.}
    The DSSYK limiting extremal surface sits at
    $\tausk=\sqrt{3/2}$ and gives
    $\CKz=2/(3\sqrt{3}\,\GN\ell^2)$: finite, covariant, $[L^{-4}]$,
    state-independent.

  \item \textbf{The double degeneracy and the gradient relation.}
    Both $\partial f/\partial\tau|_*$ and $g'(\tausk)$ vanish
    simultaneously, forcing $(\delta\tausk)^2=-3\Phi/4$ and the gradient
    relation $|\nabla\CK|=\kap\aN$.  The coupling $\kap$ cancels
    identically in the AQUAL action.

  \item \textbf{Tanh from log cosh.}
    $\mu(x)=dC_K/dx=\tanh(x)$ with correct deep-MOND and Newtonian
    limits.  The BTFR $v_c^4=\az GM_b$ follows in the deep-MOND limit.
\end{enumerate}

Testing against 175 SPARC galaxies with no free parameters, we find fits
competitive with the standard simple function.  In the most important
regime --- dwarf galaxies deep in the MOND limit --- $\mu=\tanh$
narrowly outperforms simple (median $\chi^2_r$ 7.69 vs.\ 7.88).  The
7.7\% difference at the transition is accessible to Euclid weak-lensing
stacking.

The primary open problem is deriving the AQUAL form of the kinetic action
from the DSSYK second variation.  The evidence is in hand: the
non-analytic $|\Phi|^{3/2}$ term has the correct deep-MOND power, and
$\lambda_0=7/3$ shows the MOND signal is a non-analytic family effect, not
a Jacobi instability.  Closing this step --- mapping $|\Phi|^{3/2}$ in
the DSSYK generating functional to $F(\aN^2/\az^2)$ via the holographic
dictionary --- would complete the first-principles derivation of MOND from
quantum complexity.

% ============================================================
% Figures (placeholders — generate from sparc_simulation.py)
% ============================================================

\begin{figure}[p]
  \centering
  \includegraphics[width=0.7\textwidth]{figures/fig1_mu_comparison.pdf}
  \caption{MOND interpolating functions $\mu(x)$ vs.\ $x=\aN/\az$.
    Solid blue: CODA, $\mu=\tanh(x)$, derived from
    $dC_K/dx$ with $C_K=\log\cosh$.
    Dashed green: simple, $\mu=x/\sqrt{1+x^2}$ \cite{Zhao:2006}.
    Dotted orange: standard, $\mu=x/(1+x)$ \cite{Milgrom:1983a}.
    Vertical line: MOND transition $x=1$.
    At transition: $\mu_{\tanh}(1)=0.762$ vs.\
    $\mu_{\rm simple}(1)=0.707$ --- a $+7.7\%$ difference accessible to
    Euclid weak-lensing stacking.}
  \label{fig:mu}
\end{figure}

\begin{figure}[p]
  \centering
  \includegraphics[width=0.85\textwidth]{figures/fig2_RAR.pdf}
  \caption{Radial Acceleration Relation for all 175 SPARC galaxies
    (3391 data points, grey).  Solid blue: CODA tanh,
    $\az=cH_0/2\pi=1.082\times10^{-10}\,\si{m.s^{-2}}$.
    Dashed green: tanh, $\az=\az^{\rm obs}=1.2\times10^{-10}\,\si{m.s^{-2}}$.
    Dashed and dotted white lines: deep-MOND asymptote
    $\aobs=\sqrt{\ab\az}$ and Newtonian limit $\aobs=\ab$, respectively.
    Axes normalised to $\az^{\rm obs}$.}
  \label{fig:RAR}
\end{figure}

\begin{figure}[p]
  \centering
  \includegraphics[width=\textwidth]{figures/fig3_rotation_curves.pdf}
  \caption{Rotation curves for six representative SPARC galaxies spanning
    the Hubble sequence: NGC\,3198 (Sd), NGC\,6503 (Sc), DDO\,154 (Im),
    NGC\,7331 (Sb), UGC\,2885 (Sc), IC\,2574 (Sm).
    White circles: observed $v_c\pm\sigma$.
    Solid blue: CODA tanh, $\az=cH_0/2\pi$.
    Dashed green: tanh, $\az=\az^{\rm obs}$.
    Dotted orange: simple $\mu$.
    Dash-dotted gold: baryonic Newtonian prediction.
    Inset $\chi^2_r$: CODA $|$ tanh $|$ simple.}
  \label{fig:rotation_curves}
\end{figure}

\begin{figure}[p]
  \centering
  \includegraphics[width=0.85\textwidth]{figures/fig4_delta_chi2.pdf}
  \caption{$\chi^2_r(\tanh)-\chi^2_r({\rm simple})$ per galaxy, ranked.
    Blue bars: $\tanh$ better; orange bars: simple better.
    Dashed yellow line: mean difference $+2.74$.
    The 74 galaxies where tanh outperforms simple are concentrated in
    the dwarf/deep-MOND regime.}
  \label{fig:delta_chi2}
\end{figure}

% ============================================================
\section*{Acknowledgements}

The author thanks the developers of the SPARC database (F.\ Lelli,
S.S.\ McGaugh, J.M.\ Schombert) for making rotation curve and photometry
data publicly available, and the authors of arXiv:2510.13986 and
arXiv:2403.13186 whose holographic Krylov complexity programme this work
directly extends to four dimensions.

% ============================================================
\bibliographystyle{JHEP}
\begin{thebibliography}{99}

\bibitem{Milgrom:1983a}
  M.~Milgrom,
  \emph{A modification of the Newtonian dynamics as a possible alternative to
  the hidden mass hypothesis},
  Astrophys.\ J.\ \textbf{270} (1983) 365.

\bibitem{Milgrom:1983b}
  M.~Milgrom,
  \emph{A modification of the Newtonian dynamics --- implications for galaxies},
  Astrophys.\ J.\ \textbf{270} (1983) 371.

\bibitem{Lelli:2016}
  F.~Lelli, S.S.~McGaugh and J.M.~Schombert,
  \emph{SPARC: Mass models for 175 disk galaxies with Spitzer photometry and
  accurate rotation curves},
  Astron.\ J.\ \textbf{152} (2016) 157 [arXiv:1606.09251].

\bibitem{Verheijen:2001}
  M.A.W.~Verheijen,
  \emph{The Ursa Major cluster of galaxies. V. HI rotation curve shapes and
  the Tully-Fisher relations},
  Astrophys.\ J.\ \textbf{563} (2001) 694.

\bibitem{McGaugh:2005}
  S.S.~McGaugh,
  \emph{The baryonic Tully-Fisher relation of galaxies with extended rotation
  curves and the stellar mass of rotating galaxies},
  Astrophys.\ J.\ \textbf{632} (2005) 859 [arXiv:astro-ph/0506750].

\bibitem{McGaugh:2016}
  S.S.~McGaugh, F.~Lelli and J.M.~Schombert,
  \emph{Radial acceleration relation in rotationally supported galaxies},
  Phys.\ Rev.\ Lett.\ \textbf{117} (2016) 201101 [arXiv:1609.05917].

\bibitem{Milgrom:2023}
  M.~Milgrom,
  \emph{MOND and the seven dwarfs},
  Phys.\ Rev.\ D \textbf{107} (2023) 044052.

\bibitem{Zhao:2010}
  H.S.~Zhao and B.~Li,
  \emph{Dark fluid: A unified framework for modified dynamics and dark energy},
  Astrophys.\ J.\ \textbf{712} (2010) 130.

\bibitem{Verlinde:2017}
  E.~Verlinde,
  \emph{Emergent gravity and the dark universe},
  SciPost Phys.\ \textbf{2} (2017) 016 [arXiv:1611.02269].

\bibitem{Skordis:2021}
  C.~Skordis and T.~Z{\l}o{\'s}nik,
  \emph{New relativistic theory for modified Newtonian dynamics},
  Phys.\ Rev.\ Lett.\ \textbf{127} (2021) 161302 [arXiv:2007.00082].

\bibitem{Bekenstein:2004}
  J.D.~Bekenstein,
  \emph{Relativistic gravitation theory for the modified Newtonian dynamics
  paradigm},
  Phys.\ Rev.\ D \textbf{70} (2004) 083509 [arXiv:astro-ph/0403694].

\bibitem{Deffayet:2011}
  C.~Deffayet, G.~Esposito-Far{\`e}se and R.~Woodard,
  \emph{Nonlocal metric formulations of MOND with sufficient lensing},
  Phys.\ Rev.\ D \textbf{84} (2011) 124054.

\bibitem{Parker:2019}
  D.E.~Parker, X.~Cao, A.~Avdoshkin, T.~Scaffidi and E.~Altman,
  \emph{A universal operator growth hypothesis},
  Phys.\ Rev.\ X \textbf{9} (2019) 041017 [arXiv:1812.02666].

\bibitem{Avdoshkin:2022}
  A.~Avdoshkin, A.~Dymarsky and M.~Smolkin,
  \emph{Krylov complexity in quantum field theory, and beyond},
  arXiv:2212.14429 (2022).

\bibitem{Heller:2025}
  M.P.~Heller, A.~Ori, P.~Papalini, L.~Schuhmann and Y.~Wang,
  \emph{Holographic Krylov complexity in de~Sitter space from
  double-scaled SYK},
  arXiv:2510.13986 (2025).

\bibitem{Aguilar:2024}
  S.E.~Aguilar-Gutierrez,
  \emph{Complexity in de~Sitter space from the double-scaled SYK model},
  JHEP \textbf{10} (2024) 107 [arXiv:2403.13186].

\bibitem{Zhao:2006}
  H.S.~Zhao and B.~Famaey,
  \emph{Refining the MOND interpolating function and TeVeS Lagrangian},
  Astrophys.\ J.\ \textbf{638} (2006) L9 [arXiv:astro-ph/0512425].

\bibitem{Famaey:2012}
  B.~Famaey and S.S.~McGaugh,
  \emph{Modified Newtonian dynamics (MOND): observational phenomenology
  and relativistic extensions},
  Living Rev.\ Rel.\ \textbf{15} (2012) 10 [arXiv:1112.3960].

\bibitem{Blommaert:2024}
  A.~Blommaert, T.G.~Mertens and P.~Papalini,
  \emph{DSSYK in the bulk: from Schwarzian to dilaton gravity},
  arXiv:2404.03535 (2024).

\bibitem{Brown:2016}
  A.R.~Brown, D.A.~Roberts, L.~Susskind, B.~Swingle and Y.~Zhao,
  \emph{Holographic complexity equals bulk action?},
  Phys.\ Rev.\ Lett.\ \textbf{116} (2016) 191301 [arXiv:1509.07876].

\bibitem{Riess:2022}
  A.G.~Riess et~al.,
  \emph{A comprehensive measurement of the local value of the Hubble
  constant with 1\,km/s/Mpc uncertainty from HST and the SH0ES team},
  Astrophys.\ J.\ Lett.\ \textbf{934} (2022) L7 [arXiv:2112.04510].

\bibitem{Bekenstein:1984}
  J.D.~Bekenstein and M.~Milgrom,
  \emph{Does the missing mass problem signal the breakdown of Newtonian
  gravity?},
  Astrophys.\ J.\ \textbf{286} (1984) 7.

\bibitem{Stelle:1977}
  K.S.~Stelle,
  \emph{Renormalization of higher-derivative quantum gravity},
  Phys.\ Rev.\ D \textbf{16} (1977) 953.

\bibitem{Schombert:2020}
  J.~Schombert, S.S.~McGaugh and J.~Lelli,
  \emph{Using the baryonic Tully-Fisher relation to measure $H_0$},
  Astron.\ J.\ \textbf{160} (2020) 71.

\bibitem{Amendola:2018}
  L.~Amendola et~al.\ (Euclid Theory Working Group),
  \emph{Cosmology and fundamental physics with the Euclid satellite},
  Living Rev.\ Rel.\ \textbf{21} (2018) 2 [arXiv:1606.00180].

\bibitem{Desmond:2020}
  C.~Desmond and P.G.~Ferreira,
  \emph{Galaxy morphology rules out Mond},
  Phys.\ Rev.\ D \textbf{102} (2020) 104060.

\bibitem{Gentile:2011}
  G.~Gentile, B.~Famaey and W.J.G.~de~Blok,
  \emph{THINGS about MOND},
  Astron.\ Astrophys.\ \textbf{527} (2011) A76 [arXiv:1011.4148].

\bibitem{Lelli:2017}
  F.~Lelli, S.S.~McGaugh, J.M.~Schombert and M.S.~Pawlowski,
  \emph{One law to rule them all: the radial acceleration relation of
  galaxies},
  Astrophys.\ J.\ \textbf{836} (2017) 152 [arXiv:1610.08981].

\bibitem{Angus:2006}
  G.W.~Angus, B.~Famaey and H.S.~Zhao,
  \emph{Can MOND take a bullet? Analytical comparisons of three versions of
  MOND beyond spherical symmetry},
  Mon.\ Not.\ Roy.\ Astron.\ Soc.\ \textbf{371} (2006) 138
  [arXiv:astro-ph/0606216].

\bibitem{Susskind:2022}
  L.~Susskind,
  \emph{De Sitter space, double-scaled SYK, and the separation of scales
  in the principle of relativity},
  arXiv:2209.09999 (2022).

\bibitem{CODA:dict}
  S.~Hartley,
  \emph{Holographic dictionary for CODA: from DSSYK second variation to
  AQUAL},
  CODA project, \texttt{theory/holographic\_dictionary.md} (2026).

\end{thebibliography}

\end{document}
